
\section{Proof of Theorem~\ref{thm:nested_wait_heavy_traffic_time_varying}}\label{appendix:: proof of nested optimality, time varying.}

The proof is a conservative version of the Nested WAIT argument. Instead of matching each instantaneous arrival pattern exactly, we dominate each service-window arrival count and each downstream thinning probability by worst-case quantities over the finite horizon. This gives the same bounded-queue and logarithmic-buffer structure as the constant-rate proof. The first-segment waiting condition is used only when converting bounded queues into service-normalized delay bounds; the throughput and memory conclusions use only the upper-side window-load conditions and the finite-horizon safety buffer. The structure proceeds as follows:
\begin{enumerate}
    \item \textbf{First segment analysis}: Poisson arrivals with time-varying rates are dominated by a time-invariant process using the worst-case scaled-window arrival count $\mathbf{\Lambda}^{\pi} = \sup_b \{\sum_{j=1}^m \mathcal A_j^{(\zeta)}(t(b),\Delta T_{[1,\ldots,m]})\}$. We couple this queue to a Lindley process and apply Kingman's inequality.
    \item \textbf{Subsequent segments analysis}: Binomial arrivals with time-varying thinning probabilities $p_k(b)$ are dominated by the worst-case continuation probability \(p_k^*\) defined in Section~\ref{sec:time_varying}. The same Lindley coupling then controls each downstream boundary queue.
    \item \textbf{Throughput, delay, and memory}: Bounded expected terminal work yields the normalized throughput gap. The service-normalized delay bound additionally uses the first-segment waiting condition to control idle waiting before the first segment. High-probability memory bounds follow from the same martingale construction with worst-case parameters.
\end{enumerate}

We use $b$ to denote the \emph{batch index} and $k \in \{1, \ldots, m\}$ to denote the \emph{segment index}. The total number of batches is denoted by $B$.

\subsection{First Segment Analysis}

For the first segment ($k=1$), we use the same threshold-renewal convention as in Appendix~\ref{appendix::proof of nested_wait_heavy_traffic}. Idle periods before a segment-1 service start contain fewer than \(n_1\) prompts, so they contribute only a bounded amount of terminal unfinished work and do not affect the memory bound. The first-segment waiting condition in Section~\ref{sec:time_varying} is needed later only to bound these idle periods in service-normalized time. At renewal epochs, the state transition rule is
$$ W_{(1)}^{b+1} = W_{(1)}^{b} + Y_{(1)}^b  - n_1 \cdot \mathbf{1}\{ W_{(1)}^{b} + Y_{(1)}^b \geq n_1 \},  \quad Y_{(1)}^b \sim \text{Poisson}\Big( \sum_{j=1}^m \mathcal A_j^{(\zeta)}(t(b),\Delta T_b) \Big),  $$
where $t(b)$ denotes the physical start time of batch $b$ and $\Delta T_b$ denotes the service-normalized processing time of that batch, so its physical duration is \(\Delta T_b/\zeta\).

\paragraph{Worst-case domination.} Time-varying arrival rates prevent exact load balance because the arrival composition changes over the horizon. We therefore use a conservative domination argument. Since $\Delta T_b\leq \Delta T_{[1,\dots,m]}$ (the service-normalized time for a full batch containing all $m$ types), the arrivals during any realized service period are dominated by arrivals over a full service-normalized window. This gives the process $\{\bar{W}_{(1)}^{b}\}$:
$$ \bar{W}_{(1)}^{b+1} = \bar{W}_{(1)}^{b} + \bar{Y}_{(1)}^b  - n_1 \cdot \mathbf{1}\{ \bar{W}_{(1)}^{b} + \bar{Y}_{(1)}^b \geq n_1 \},  \quad \bar{Y}_{(1)}^b \sim \text{Poisson}\Big(  \sum_{j=1}^m \mathcal A_j^{(\zeta)}(t(b),\Delta T_{[1,\dots,m]}) \Big).  $$
By~(\ref{eq:nested_wait_thresholds_time_varying}), the aggregate arrival rate during any batch is bounded by the worst-case aggregate rate:
$$ \sum_{j=1}^m \mathcal A_j^{(\zeta)}(t(b),\Delta T_{[1,\dots,m]}) \leq \sup_{\forall b}\left\{\sum_{j=1}^m \mathcal A_j^{(\zeta)}(t(b),\Delta T_{[1,\dots,m]}) \right\} < n_1.$$
Define $\mathbf{\Lambda}^{\pi} = \sup_{\forall b}\left\{\sum_{j=1}^m \mathcal A_j^{(\zeta)}(t(b),\Delta T_{[1,\dots,m]}) \right\} < n_1$ as the worst-case aggregate arrival count over a full service-normalized iteration window. This enables a time-invariant dominating process that provides tractable bounds:
$$ \tilde{W}_{(1)}^{b+1} = \tilde{W}_{(1)}^{b} + \tilde{Y}_{(1)}^b  - n_1 \cdot \mathbf{1}\{ \tilde{W}_{(1)}^{b} + \tilde{Y}_{(1)}^b \geq n_1 \},  \quad \tilde{Y}_{(1)}^b \sim \text{Poisson}(\mathbf{\Lambda}^{\pi}).  $$
This time-invariant process enables the same queueing bound used in the constant-rate proof. Define $X_{(1)}^b = \tilde{Y}_{(1)}^b - n_1$. The process $\{ \tilde{W}_{(1)}^{b}\}$ is dominated by a Lindley process:
$$ \overline{W}_{(1)}^{0} = 2n_1,\quad  \overline{W}_{(1)}^{b+1} = \max \{  2n_1, \overline{W}_{(1)}^{b} + X_{(1)}^b \},\quad \forall b \in \mathbb{N}. $$
Thus we have the dominance chain:
$$ W_{(1)}^{b}\leq \bar{W}_{(1)}^{b} \leq \tilde{W}_{(1)}^{b} \leq \overline{W}_{(1)}^{b}, \quad \forall b \in \mathbb{N}.$$
Since $n_1> \mathbf{\Lambda}^{\pi}$, the process $\{\overline{W}_{(1)}^{b} \}$ has negative drift. By Kingman's inequality \citep{kingman1962queues}, with
$$\text{Var}(X_{(1)}^b ) = \mathbf{\Lambda}^{\pi},\quad \left|  \ex{}{X_{(1)}^b} \right| = n_1 - \mathbf{\Lambda}^{\pi},$$
we obtain
$$  \ex{}{W_{(1)}^b}\leq  2n_1 + \frac{\mathbb{E}[(\tilde{Y}_{(1)}^b-n_1)^2]}{2\left(n_1 -\mathbf{\Lambda}^{\pi}\right)}. $$
\subsection{Subsequent Segments Analysis ($k \geq 2$)}

For segment $k$ ($2\leq k\leq m$), arrivals come from prompts that completed segment $k-1$. The batch index $b$ for segment $k$ counts only batches containing segment $k-1$. These indices differ across segments but are bounded by the total batch count $B$. As in Appendix~\ref{appendix::proof of nested_wait_heavy_traffic}, \(W^b_{(k)}\) is the post-review carryover residual at the downstream boundary; fresh survivors from segment \(k-1\) are charged to the upstream threshold batch, and only residual carryover receives the extra safety buffer.

\paragraph{Time-varying thinning probabilities.} Consider the $b$-th arrival to segment $k$, which comes from the $b$-th batch containing segment $k-1$. The thinning probability $p_k(b)$ depends on the prompt composition in this batch and varies over time. Although exact arrival times are unknown, the relevant prompts accumulate over one or more scaled windows \((t,h)\) before the batch containing segment \(k-1\) begins. To handle this time variation, we use the worst-case thinning probability. Conditional on the filtration before boundary revelation, the survivor indicators in a boundary batch are independent Bernoulli random variables with possibly different probabilities \(q_i\le p_k^*\). Their sum is therefore a Poisson-binomial random variable stochastically dominated by \(\mathrm{Binomial}(n_{k-1},p_k^*)\). By~(\ref{eq:nested_wait_thresholds_time_varying}) and the definition of \(p_k^*\) in Section~\ref{sec:time_varying},
$$p_k(b) \leq p_k^* < \frac{n_k}{n_{k-1}}, \quad \forall b \in \mathbb{N}.$$
The arrival process is
$$Y_{(k)}^b \sim \text{Binomial}(n_{k-1},p_k(b)),\quad \forall b\geq 0,$$
and the state transition rule is
$$W_{(k)}^{b+1} = W_{(k)}^{b} + Y_{(k)}^{b} - n_k \cdot\mathbf{1}\{ W_{(k)}^{b} + Y_{(k)}^{b}\geq n_k \}, $$
where $b$ indexes batches containing segment $k-1$, since new arrivals to segment $k$ occur only when segment $k-1$ is processed.

Define $X_{(k)}^b = Y_{(k)}^b - n_k$. Although the distribution $Y_{(k)}^b\sim \text{Binomial}(n_{k-1},p_k(b))$ is time-varying, the coupling construction follows the same induction as Appendix~\ref{appendix_lemma::proof of nested coupled process}:
\begin{lemma}[Coupling for Time-Varying Case]\label{lemma: coupling, time varying}
Define a coupled process $\{ \tilde{W}_{(k)}^b\}$ by:
\begin{equation*}
\begin{aligned}
\tilde{W}_{(k)}^0 &= n_k, \\
\tilde{W}_{(k)}^{b+1} &= \max\{ n_k, \, \tilde{W}_{(k)}^b+X_{(k)}^b\}, \\
X_{(k)}^b &= Y_{(k)}^b - n_k.
\end{aligned}
\end{equation*}
Then $\tilde{W}^b_{(k)}\geq W^b_{(k)}$ for all $b\in \mathbb{N}$.
\end{lemma}
To apply Kingman's inequality \citep{kingman1962queues}, we construct a time-invariant dominating process:
\begin{equation*}
\begin{aligned}
\bar{W}_{(k)}^0 &= n_k, \\
\bar{W}_{(k)}^{b+1} &= \max\{ n_k,\, \bar{W}_{(k)}^{b}+\bar{X}_{(k)}^b \}, \\
\bar{X}_{(k)}^b &= \bar{Y}_{(k)}^b  -n_k, \\
\bar{Y}_{(k)}^b &\sim \text{Binomial}(n_{k-1}, p_k^*), \quad \forall b\in \mathbb{N}.
\end{aligned}
\end{equation*}
Then we have the dominance chain:
$$\bar{W}^b_{(k)}\geq \tilde{W}^b_{(k)}\geq W^b_{(k)},\quad \forall b\in \mathbb{N}.$$
The process $\bar{W}^b_{(k)}$ has the Lindley representation:
\begin{equation*}
\begin{aligned}
\bar{W}_{(k)}^0 &= n_k, \\
\bar{W}_{(k)}^{b+1} &= n_k + \max_{0\leq i\leq b}\{ \bar{S}_{(k)}^i \}, \\
\text{where } \bar{S}_{(k)}^0 &= 0,\; \bar{S}_{(k)}^i = \sum_{r=1}^i \bar{X}_{(k)}^{r},\; 1\leq i\leq b.
\end{aligned}
\end{equation*}
Since
$$ \ex{}{\bar{X}_{(k)}^b} = n_{k-1}\cdot p_{k}^* - n_k <0,\quad \forall b\in \mathbb{N}, $$
the coupled process $\{\bar{W}_{(k)}^b \}$ has negative drift. With
$$\text{Var}(\bar{X}_{(k)}^b) = n_{k-1}p_k^*(1-p_k^*),\quad \left|\ex{}{\bar{X}_{(k)}^b}\right|=n_k - p_k^* n_{k-1},$$
Kingman's inequality \citep{kingman1962queues} yields:
$$\ex{}{\max_{0\leq i\leq b}\{ \bar{S}_{(k)}^i \}  }\leq  \frac{  n_{k-1}p_k^*(1-p_k^*) }{ 2(n_k - n_{k-1}p_k^*)}. $$
Thus, the expected queue length is bounded:
$$\ex{}{W_{(k)}^b}\leq \ex{}{\tilde{W}_{(k)}^b}\leq \ex{}{\bar{W}_{(k)}^b} \leq n_k + \frac{\mathbb{E}[(\bar{Y}_{(k)}^b-n_k)^2]}{2(n_k - n_{k-1}p_k^*)}.$$
\subsection{Throughput, Delay, and High-Probability Bounds}

The bounded queue estimates above imply that the expected terminal unfinished work is \(O(1)\). This includes the first-segment entry queue, the downstream boundary queues, the bounded segment interiors, and the final partial service opportunity. Dividing this terminal work by the scaled horizon \(\zeta T\) gives
\[
    \throughput_T^*-\mathbb{E}[\throughput^{(\zeta,\pi)}]
    =
    O((\zeta T)^{-1}).
\]
This argument does not require a lower bound on the arrival rate. If arrivals fall below the first threshold near the end of the horizon, fewer than \(n_1\) prompts can remain in the external entry queue, which is still a bounded terminal effect.

For service-normalized delay, the queue-integral argument from Appendix~\ref{appendix::proof of nested_wait_heavy_traffic} applies to the service periods and downstream boundary queues because their expected lengths are uniformly bounded. The only additional issue in the time-varying setting is waiting for the Segment 1 entry queue to reach its threshold during low-arrival intervals. The first-segment waiting condition gives \(\sup_t\zeta\,\mathbb E[\tau_1^{(\zeta)}(t)]<\infty\), so these idle waiting periods add only \(O(1)\) in service-normalized time. Hence
\[
    \mathbb{E}[\Latency^{(\zeta,\pi)}],
    \mathbb{E}[\TTFT^{(\zeta,\pi)}]
    =
    O(1)
\]
whenever the first-segment waiting condition holds.

For high-probability memory bounds, the analysis follows the same martingale construction applied to the time-invariant coupled process $\{ \bar{W}_{(k)}^b \}$ defined above. Each realized iteration has duration at least \(d_0/\zeta\), so the number of realized opportunities over \([0,T]\) is at most \(B_\zeta^{\max}=\lceil \zeta T/d_0\rceil\). By the same argument as in Appendix~\ref{appendix::proof of nested_wait_heavy_traffic}, with $\bar{X}_{(k)}^{r} = \bar{Y}_{(k)}^r - n_k$, $\bar{Y}_{(k)}^r \sim \text{Binomial}(n_{k-1}, p_k^*)$, and $p_k^*<n_k/n_{k-1}<1$, the memory bound is:
$$\mathbf{M}^{(\zeta,T)} = M^{\pi} + \sum_{k=2}^m n_k \cdot (l + l_{k-1}') + \sum_{k=2}^m \theta_k^{-1}\ln\left( \frac{m(1+\zeta T/d_0)}{\delta} \right) \cdot (l + l_{k-1}'),$$
where $\theta_k$ ($2\leq k\leq m$) is the unique positive solution to:
\begin{equation}
e^{-\theta_k n_k} (1 - p_k^* + p_k^* e^{\theta_k})^{n_{k-1}} = 1.\label{eq::eq of theta, time varying}
\end{equation}
Since \(p_k^*\) is used to dominate every realized thinning probability, the same exponential-martingale argument as Appendix~\ref{appendix_lemma::proof of solution to martingale exists} applies with this worst-case binomial increment. The lower bound for \(\theta_k\) follows from the same calculation.
